% =====================================================================
%  Formai beállítások a „A szakdolgozat formai és tartalmi követelményei”
%  dokumentum alapján (a pontszámok arra utalnak)
% =====================================================================

% --- Betűtípus (24. pont: Times New Roman) – XeLaTeX kell hozzá -------
\usepackage{fontspec}
\setmainfont{Times New Roman}
\setmonofont{Consolas}[Scale=0.88]   % programkódhoz

% --- Nyelv ------------------------------------------------------------
\usepackage[english,magyar]{babel}   % az utolsó a fő nyelv
\usepackage{csquotes}                 % \enquote{...} -> „...”

% --- Margók (25. pont: bal 3 cm, többi 2,5 cm), A4 --------------------
\usepackage[a4paper,left=3cm,right=2.5cm,top=2.5cm,bottom=2.5cm]{geometry}

% --- Sortávolság (25. pont: másfeles) ---------------------------------
% A 12pt-s report osztályban a lábjegyzet (\footnotesize) eleve 10pt.
\usepackage{setspace}
\onehalfspacing

% --- Címsorok (26. pont: fejezetcím középre, a többi balra) -----------
% (27. pont: minden fejezet új oldalon kezdődik – ezt a report osztály adja)
\usepackage{titlesec}
\titleformat{\chapter}{\normalfont\Large\bfseries\centering}{\thechapter.}{0.6em}{}
\titleformat{name=\chapter,numberless}{\normalfont\Large\bfseries\centering}{}{0pt}{}
\titlespacing*{\chapter}{0pt}{0pt}{1.5em}
\titleformat{\section}{\normalfont\large\bfseries}{\thesection.}{0.6em}{}
\titleformat{\subsection}{\normalfont\normalsize\bfseries}{\thesubsection.}{0.6em}{}
\setcounter{secnumdepth}{2}
\setcounter{tocdepth}{2}

% --- Ábrák, táblázatok (9–10. pont: számozva, címmel) -----------------
\usepackage{graphicx}
\graphicspath{{abrak/}}
\usepackage{booktabs}
\usepackage{tabularx}
\usepackage{array}
\usepackage[font=small,labelfont=bf]{caption}
\usepackage{float}
\usepackage{enumitem}
\setlist{nosep}

% --- Programkód -------------------------------------------------------
\usepackage{xcolor}
\usepackage{listings}
\renewcommand{\lstlistingname}{kódrészlet}
\renewcommand{\lstlistlistingname}{Kódrészletek jegyzéke}
\lstset{
  basicstyle=\ttfamily\small\setstretch{1},
  keywordstyle=\color{blue!60!black}\bfseries,
  commentstyle=\color{green!40!black}\itshape,
  stringstyle=\color{red!50!black},
  numbers=left, numberstyle=\tiny\color{gray}, numbersep=6pt,
  frame=single, rulecolor=\color{gray!50},
  breaklines=true, showstringspaces=false, tabsize=4,
  captionpos=b, columns=fullflexible, keepspaces=true,
}
\lstdefinelanguage{TypeScript}{
  language=Java,
  morekeywords={const,let,async,await,function,interface,type,export,import,from,useState,useEffect},
  morestring=[b]', morestring=[b]`,
}

% --- Irodalomjegyzék (7., 8., 11. pont) -------------------------------
% Egységes, számozott hivatkozás: [3], szó szerinti idézetnél oldalszámmal: [3, 45. o.]
\usepackage[backend=biber,style=numeric,sorting=none,urldate=iso,seconds=true]{biblatex}
\addbibresource{irodalom.bib}
\DefineBibliographyStrings{magyar}{urlseen={Letöltve:}}

% --- Mellékletek, PDF lapok beillesztése (feladatkiírás stb.) ---------
\usepackage{pdfpages}

% --- Kattintható hivatkozások (utolsónak) -----------------------------
\usepackage[hidelinks]{hyperref}

% --- Saját parancsok --------------------------------------------------
% Kitöltendő helyek jelölése: pirosan látszik a PDF-ben, hogy ne maradjon benne
\newcommand{\todo}[1]{\textcolor{red}{\textbf{[#1]}}}
% Ha a megadott PDF létezik, beilleszti; különben egy helyőrző oldalt tesz be
\newcommand{\beillesztPDF}[2]{%
  \IfFileExists{#1}{\includepdf[pages=-]{#1}}{%
    \thispagestyle{empty}\null\vfill
    \begin{center}\Large\textcolor{red}{\textbf{[#2]}}\\[1em]
    \normalsize Ide kerül az aláírt, beszkennelt lap: \texttt{\detokenize{#1}}\end{center}\vfill\clearpage}}
